% TOP_PROBLEM: {"id": 30006727, "title": "L versus NL", "category_id": 15, "category": {"id": 15, "name": "computer_science", "display_name": "Computer Science", "description": "Computational complexity, algorithms, and theoretical CS.", "slug": "computer-science", "order_index": 15, "created_at": "2026-07-31T15:26:25.671Z"}, "status": "open"}
% FIRST_POSTED: 2026-09-25
% ENTRY_KIND: statement_only
% !TeX program = lualatex
% Definition and sources only; this file is not an LLM solution attempt.
\documentclass[11pt]{article}
\usepackage[margin=25mm]{geometry}
\usepackage{fontspec}
\setmainfont{Latin Modern Roman}
\usepackage{amsmath,amssymb,mathtools}
\usepackage{unicode-math}
\setmathfont{Latin Modern Math}
\usepackage{xurl,hyperref}
\hypersetup{colorlinks=true,urlcolor=blue}
\setcounter{secnumdepth}{0}
\setlength{\parindent}{0pt}
\setlength{\parskip}{5pt}
\setlength{\emergencystretch}{3em}
\begin{document}
\section*{\texorpdfstring{L versus NL}{L versus NL}}\label{30006727}
\subsection{Definitions and mathematical statement}
The classes $\mathsf L$ and $\mathsf{NL}$ use deterministic and nondeterministic $O(\log n)$ work space, respectively, with a read-only input tape. In nondeterministic acceptance, a yes-input has an accepting branch and a no-input has none. The equivalent graph problem asks for a deterministic logarithmic-space algorithm deciding whether
\[
 \exists\text{ directed path }s\leadsto t\text{ in }G.
\]
The graph is arbitrary and directed. Undirected connectivity, extra work space, or a nonuniform circuit procedure does not settle this target.
\par\smallskip\textit{Formulation and concept references:} \href{https://drops.dagstuhl.de/storage/00lipics/lipics-vol099-socg2018/LIPIcs.SoCG.2018/LIPIcs.SoCG.2018.pdf}{[S1]}.

\subsection{Short English statement}
Can nondeterministic logarithmic-space computation always be simulated with the same asymptotic amount of deterministic working memory?
\subsection{Sources}
\begin{itemize}
\item[S1]Space-efficient algorithms for reachability in directed geometric graphs.
\url{https://drops.dagstuhl.de/storage/00lipics/lipics-vol099-socg2018/LIPIcs.SoCG.2018/LIPIcs.SoCG.2018.pdf}.
\textit{Formulation source.}
\end{itemize}
\end{document}
